\clearpage\SourceLeaf{302}
\begin{center}
\textbf{[ON MÖBIUS'S BARYCENTRIC CALCULUS.]}

\medskip

[I.]
\end{center}

\begin{center}\textsc{Gauss} to \textsc{Schumacher}. Göttingen, 12 May 1843.\end{center}

I still owe you my thanks, my dearest friend, for several of your letters and other
interesting communications. But you must not keep too strict an account with
me. For 4---6 weeks I have been drawn into some
mathematical speculations (at first by chance circumstances), in which new prospects
have repeatedly led me in different directions. I have achieved much and missed much. Amid
such wandering one is (or at least I am) almost incapable of taking up any
other subject, and so my correspondence in particular has fallen very far behind.
Those speculations have mostly concerned carrying through methods of my own
rather than new matters. Most recently they have concerned, among other things, several
questions about conics. This has reminded me that half a
century ago, when I first read \textsc{Newton}'s Principia, I found several things
unsatisfactory, especially his propositions on conics, admirable in themselves.
Yet I always read with the feeling that what I had learned did not give me command of the
subject. I was particularly troubled by the straight line by means of which a conic
can be described. \textsc{Newton} solves the problem of drawing a conic through 5 points $A$, $B$,
$C$, $D$, $E$ by first using the four
